\documentclass[11pt]{article}
\usepackage[T1]{fontenc}
\usepackage{amsmath,amssymb,amsthm,lmodern}
\usepackage[margin=1in]{geometry}
\usepackage{hyperref}
\hypersetup{hidelinks,pdftitle={EGNO: The coend construction, Theorem 1.10.1},pdfauthor={Pavel Etingof; Shlomo Gelaki; Dmitri Nikshych; Victor Ostrik}}
\newcommand{\kk}{\Bbbk}
\newcommand{\cC}{\mathcal{C}}
\newcommand{\cD}{\mathcal{D}}
\newcommand{\FiniteVec}{\mathrm{Vec}}
\newcommand{\AllVec}{\mathbf{Vec}}
\newcommand{\Coend}{\operatorname{Coend}}
\newcommand{\End}{\operatorname{End}}
\newcommand{\Hom}{\operatorname{Hom}}
\newcommand{\Ker}{\operatorname{Ker}}
\newcommand{\Ind}{\operatorname{Ind}}
\newcommand{\id}{\operatorname{id}}
\newcommand{\comod}{\mathrm{comod}}
\newcommand{\modu}{\mathrm{mod}}
\newcommand{\Rad}{\operatorname{Rad}}
\newcommand{\Spec}{\operatorname{Spec}}
\newcommand{\sourceloc}[1]{\par\smallskip\noindent\textit{Source: #1.}\par\smallskip}
\newcommand{\editorial}[1]{\par\smallskip\noindent\textsf{Editorial note.} #1\par\smallskip}
\title{The Coend Construction\\Theorem 1.10.1}
\author{Pavel Etingof \and Shlomo Gelaki \and Dmitri Nikshych \and Victor Ostrik}
\date{Private reading transcription of the 2015 author final copy\\with a separately identified author clarification dated February 14, 2026}
\begin{document}
\maketitle

% BEGIN CONTEXT EGNO:1.10.1
\noindent\textbf{Source and scope.}
\emph{Tensor Categories}, Mathematical Surveys and Monographs, volume 205,
American Mathematical Society, 2015, author final copy (362 physical pages).
The fixed source is
\href{https://math.mit.edu/~etingof/egnobookfinal.pdf}{\nolinkurl{https://math.mit.edu/~etingof/egnobookfinal.pdf}};
the authors' book page is
\href{https://math.mit.edu/~etingof/books.html}{\nolinkurl{https://math.mit.edu/~etingof/books.html}}.
The main selection is the complete Section 1.10 on physical page 31
(printed page 14), together with the statement of Theorem 1.9.15 on that page,
its complete proof and footnote 4 on physical page 32 (printed page 15),
and Exercise 1.10.2(i) on physical page 32.
Theorem 1.9.15 and the checking exercise are supporting material for this
single primary target, Theorem 1.10.1.

The separate six-page document
\emph{Corrections to the book ``Tensor categories'' by Etingof, Gelaki,
Nikshych and Ostrik, AMS, 2015}, dated February 14, 2026, is available at
\href{https://math.mit.edu/~etingof/tensbookcor.pdf}{\nolinkurl{https://math.mit.edu/~etingof/tensbookcor.pdf}}.
Its clarification of Subsection 1.10, on its physical page 1, is reproduced
separately below; it has not been silently substituted into the 2015 text.

\editorial{This is a private source transcription. The authors explicitly call
the primary proof a sketch; its three verification statements remain at the
authors' stated level of detail. Basic category theory, ind-completions,
injective hulls, and Grothendieck categories are background. No new proof of
the authors' assertions is supplied. Exercise 1.10.2(ii), concerning
$\mathcal O(G)$, is a separate example and is not part of this selection.
Original bracketed reference keys are retained for source identification.}

\noindent\textbf{Original source footer.}
``Author's final version made available with permission of the publisher,
American Mathematical Society. License or copyright restrictions may apply
to redistribution; see
\url{http://www.ams.org/publications/ebooks/terms}.''
The author page restricts posting this manuscript on other websites; this
reading copy is kept private.

\subsection*{Conventions and terminology from the source}
\sourceloc{Section 1.1, physical page 18, printed page 1; selected conventions}
The identity endofunctor of a category $\cC$ will be denoted $\id_{\cC}$.
We say that a functor $F:\cC\to\cD$ is an \emph{equivalence} if there is a
functor $F^{-1}:\cD\to\cC$, called a \emph{quasi-inverse} of $F$, such that
$F^{-1}\circ F\cong\id_{\cC}$ and $F\circ F^{-1}\cong\id_{\cD}$.

A category is called \emph{locally small} if for any objects $X,Y$,
$\Hom_{\cC}(X,Y)$ is a set, and is called \emph{essentially small} if in
addition its isomorphism classes of objects form a set. In other words,
an essentially small category is a category equivalent to a small category
(i.e., to one where objects and morphisms form a set). All categories
considered in this book will be locally small (except in the section on
2-categories), and most of them will be essentially small. So set-theoretical
subtleties will not play any role in this book.\footnote[1]{A well-known
foundational subtlety of category theory (which, luckily, does not play an
essential role in most category-theoretical considerations) is that in
general, objects and morphisms of a category are defined to form a class
rather than a set (this allows one to consider the category of all sets
without running into Russell's paradox---the non-existence of the set of
all sets). For example, the category of all sets or all vector spaces over
a field is not small or even essentially small (although it is locally
small). However, in this book we will usually consider categories with
finiteness conditions, such as the category of finite dimensional vector
spaces, so this issue is not going to be relevant most of the time. In fact,
even when we work with categories that are not essentially small (such as
the category of all vector spaces), we will allow ourselves to abuse
terminology and speak about ``the set of isomorphism classes of objects''
of such a category. This does not create a confusion in our context.}

For a category $\cC$ the notation $X\in\cC$ will mean that $X$ is an object
of $\cC$, and the set of morphisms between $X,Y\in\cC$ will be denoted by
$\Hom_{\cC}(X,Y)$. An element $\phi\in\Hom_{\cC}(X,Y)$ will be usually
depicted either as $\phi:X\to Y$ or as $X\xrightarrow{\phi}Y$.
We denote by $\cC^{\vee}$ the category \emph{dual} to $\cC$, i.e., the one
obtained from $\cC$ by reversing the direction of morphisms.

Unless otherwise specified, all fields are assumed to be algebraically closed.

\sourceloc{Example 1.3.3, physical page 20, printed page 3}
The category of abelian groups is an abelian category. The category of modules
over a ring is an abelian category. The category $\AllVec$ of vector spaces
over a field $\kk$ and its subcategory $\FiniteVec$ of finite dimensional vector
spaces are $\kk$-linear abelian categories. More generally, the category of
modules over an associative $\kk$-algebra and the category of comodules over
a coassociative $\kk$-coalgebra (see Section 1.9 below) are $\kk$-linear
abelian categories.

\editorial{The bold $\AllVec$ and ordinary $\FiniteVec$ reproduce the source's
distinction between all vector spaces and finite dimensional vector spaces.
The functors to $\FiniteVec$ below use the latter. The general convention about
fields remains in force. Essential smallness is stated explicitly in
Theorem 1.9.15; it is not added to the wording of Theorem 1.10.1.}

\sourceloc{Definitions 1.8.1 and 1.8.5, physical page 26, printed page 9}
Let $\kk$ be a field.

\noindent\textbf{Definition 1.8.1.}
A $\kk$-linear abelian category $\cC$ is said to be \emph{locally finite}
if the following two conditions are satisfied:
\begin{enumerate}
\item[(i)] for any two objects $X,Y$ in $\cC$ the $\kk$-vector space
  $\Hom_{\cC}(X,Y)$ is finite dimensional;
\item[(ii)] every object in $\cC$ has finite length.
\end{enumerate}

\noindent\textbf{Definition 1.8.5.}
A $\kk$-linear abelian category $\cC$ is said to be \emph{finite} if it is
equivalent to the category $A$-$\modu$ of finite dimensional modules over
a finite dimensional $\kk$-algebra $A$.

\sourceloc{Definitions 1.9.1 and 1.9.2, physical page 29, printed page 12}
\noindent\textbf{Definition 1.9.1.}
A \emph{coalgebra} (with counit) over a field $\kk$ is a $\kk$-vector space
$C$ together with a \emph{comultiplication} (or \emph{coproduct})
$\Delta:C\to C\otimes C$ and \emph{counit} $\varepsilon:C\to\kk$ such
that $\Delta$ and $\varepsilon$ are $\kk$-linear maps and
\begin{enumerate}
\item[(i)] $\Delta$ is coassociative, i.e.,
  \[
  (\Delta\otimes\id_C)\circ\Delta
  =(\id_C\otimes\Delta)\circ\Delta
  \]
  as maps $C\to C^{\otimes 3}$;
\item[(ii)] one has
  \[
  (\varepsilon\otimes\id_C)\circ\Delta
  =(\id_C\otimes\varepsilon)\circ\Delta=\id_C
  \]
  as maps $C\to C$ (the ``counit axiom'').
\end{enumerate}

\noindent\textbf{Definition 1.9.2.}
A \emph{left comodule} over a coalgebra $C$ is a vector space $M$ together
with a linear map $\pi:M\to C\otimes M$ (called the coaction map), such
that for any $m\in M$, one has
\[
(\Delta\otimes\id)(\pi(m))=(\id\otimes\pi)(\pi(m)),
\qquad (\varepsilon\otimes\id)(\pi(m))=m.
\]
Similarly, a \emph{right comodule} over $C$ is a vector space $M$ together
with a linear map $\pi:M\to M\otimes C$, such that for any $m\in M$, one has
\[
(\pi\otimes\id)(\pi(m))=(\id\otimes\Delta)(\pi(m)),
\qquad (\id\otimes\varepsilon)(\pi(m))=m.
\]
For example, $C$ is a left and right comodule over itself with $\pi=\Delta$.

\sourceloc{Definition 1.9.13 and Remark 1.9.14, physical page 30, printed page 13}
\noindent\textbf{Definition 1.9.13.}
A coalgebra $C$ is \emph{pointed} if any simple right $C$-comodule is
1-dimensional.

\noindent\textbf{Remark 1.9.14.}
A finite dimensional coalgebra $C$ is pointed if and only if the algebra
$C^*$ is \emph{basic}, i.e., the quotient $C^*/\Rad(C^*)$ of $C^*$ by its
radical is commutative. In this case, simple $C$-comodules are points of
$\Spec(C^*/\Rad(C^*))$ (here $\Spec$ stands for the set of maximal ideals),
which justifies the term ``pointed''.

% BEGIN SUPPORT STATEMENT EGNO:1.9.15
\sourceloc{Theorem 1.9.15, physical page 31, printed page 14}
In the next section we will prove the following theorem:

\noindent\textbf{Theorem 1.9.15 (Takeuchi, [Tak2]).}
\emph{Any essentially small locally finite abelian category $\cC$ over a
field $\kk$ is equivalent to the category $C$-$\comod$ for a unique pointed
coalgebra $C$. In particular, if $\cC$ is finite, it is equivalent to the
category $A$-$\modu$ for a unique basic algebra $A$ (namely, $A=C^*$).}
% END SUPPORT STATEMENT EGNO:1.9.15

\subsection*{1.10. The Coend construction}
\sourceloc{physical page 31, printed page 14; complete preamble of Section 1.10}
Let $\cC$ be a $\kk$-linear abelian category, and $F:\cC\to\FiniteVec$ an exact,
faithful functor. In this case one can define the space $\Coend(F)$ as follows:
\begin{equation}
\Coend(F):=\left(\bigoplus_{X\in\cC}F(X)^*\otimes F(X)\right)/E
\tag{1.9}\label{eq:original-1.9}
\end{equation}
where $E$ is spanned by elements of the form
$y_*\otimes F(f)x-F(f)^*y_*\otimes x$,
$x\in F(X)$, $y_*\in F(Y)^*$, $f\in\Hom(X,Y)$; in other words,
\[
\Coend(F)=\varinjlim\End(F(X))^*.
\]
Thus we have $\End(F)=\varprojlim\End(F(X))=\Coend(F)^*$, which yields
a coalgebra structure on $\Coend(F)$. So the algebra $\End(F)$ (which may
be infinite dimensional) carries the inverse limit topology, in which
a basis of neighborhoods of zero is formed by the kernels $K_X$ of the maps
$\End(F)\to\End(F(X))$, $X\in\cC$, and
$\Coend(F)=\End(F)^{\vee}$, the space of continuous linear functionals
on $\End(F)$.

The following theorem is standard (see [Tak2]).

% END CONTEXT EGNO:1.10.1

% BEGIN PRIMARY STATEMENT EGNO:1.10.1
\sourceloc{Theorem 1.10.1, physical page 31, printed page 14}
\noindent\textbf{Theorem 1.10.1.}
\emph{Let $\cC$ be a $\kk$-linear abelian category with an exact faithful
functor $F:\cC\to\FiniteVec$. Then $\cC$ is locally finite, and $F$ defines an
equivalence between $\cC$ and the category of finite dimensional right
comodules over $C:=\Coend(F)$ (or, equivalently, with the category of
continuous finite dimensional left $\End(F)$-modules).}
% END PRIMARY STATEMENT EGNO:1.10.1

% BEGIN PRIMARY PROOF EGNO:1.10.1
\sourceloc{physical page 31, printed page 14; complete author proof sketch}
\begin{proof}
We sketch the proof, leaving the details to the reader. Consider the
ind-object $Q:=\bigoplus_{X\in\cC}F(X)^*\otimes X$.
For $X,Y\in\cC$ and $f\in\Hom(X,Y)$, let
\[
j_f:F(Y)^*\otimes X\longrightarrow
F(X)^*\otimes X\oplus F(Y)^*\otimes Y\subset Q
\]
be the morphism defined by the formula
\[
j_f=\id\otimes f-F(f)^*\otimes\id.
\]
Let $I$ be the quotient of $Q$ by the image of the direct sum of all $j_f$.

The following statements are not hard to verify:
\begin{enumerate}
\item[(i)] $I$ represents the functor $F(-)^*$, i.e., $\Hom(X,I)$ is
  naturally isomorphic to $F(X)^*$; in particular, $I$ is injective.
\item[(ii)] $F(I)=C$, and $I$ is naturally a left $C$-comodule. Its comodule
  structure is induced by the coevaluation morphism
  \[
  F(X)^*\otimes X
  \xrightarrow{\operatorname{coev}_{F(X)}}
  F(X)^*\otimes F(X)\otimes F(X)^*\otimes X.
  \]
\item[(iii)] Let us regard $F$ as a functor $\cC\to C$-$\comod$.
  For $M\in C$-$\comod$, let
  $\theta_M:M\otimes I\to M\otimes C\otimes I$ be the morphism
  $\pi_M\otimes\id-\id\otimes\pi_I$, and let $K_M$ be the kernel of
  $\theta_M$. Then the functor $G:C$-$\comod\to\cC$ given by the formula
  $G(M)=\Ker\theta_M$, is a quasi-inverse to $F$.
\end{enumerate}
This completes the proof.
\end{proof}
% END PRIMARY PROOF EGNO:1.10.1

% BEGIN SUPPORT EXERCISE EGNO:1.10.2(i)
\sourceloc{Exercise 1.10.2(i), physical page 32, printed page 15}
\noindent\textbf{Exercise 1.10.2.}
(i) Check statements (i)--(iii) in the proof of Theorem 1.10.1.
% END SUPPORT EXERCISE EGNO:1.10.2(i)

% BEGIN SUPPORT PROOF EGNO:1.9.15
\sourceloc{physical page 32, printed page 15; complete proof of Theorem 1.9.15,
ending before Section 1.11}
\noindent\textbf{Proof of Theorem 1.9.15.}
Now let us use the Coend construction to prove Theorem 1.9.15.
Consider the dual category $\cC^{\vee}$ (which is also locally finite).
It is known that the ind-completion of $\cC^{\vee}$ is a Grothendieck
category, so it has enough injectives.\footnote[4]{The notion of a
Grothendieck category comes from the paper [Gr]; see [KashS] for basics
of the theory of Grothendieck categories.}
This means that every simple object $L$ of $\cC^{\vee}$ has an injective
hull $Q_L$ in the ind-completion $\Ind(\cC^{\vee})$.
Let $Q=\bigoplus_L Q_L$, where the sum is taken over isomorphism classes
of simple objects of $\cC^{\vee}$. Consider the (covariant) functor
$F:\cC\to\FiniteVec$ given by $F(X)=\Hom_{\cC^{\vee}}(X,Q)$.
The functor $F$ is exact and faithful, and for every simple object
$L\in\cC$, $\dim F(L)=1$. This means that $C:=\Coend(F)$ is a pointed
coalgebra, and by Theorem 1.10.1 we have an equivalence of categories
$C$-$\comod\cong\cC$. Moreover, suppose $C'$ is another pointed
coalgebra, and $G:C'$-$\comod\to C$-$\comod$ is an equivalence.
Then $Q=G(C')$ is an injective object in the ind-completion of
$C$-$\comod^{\vee}$ such that $\dim\Hom(L,Q)=1$ for any simple object $L$.
This implies that $Q=C$, hence $C'$ is isomorphic to $C$.
Theorem 1.9.15 is proved.
% END SUPPORT PROOF EGNO:1.9.15

% BEGIN SUPPORT AUTHOR CLARIFICATION EGNO:1.10.1
\clearpage
\subsection*{Separate author clarification: February 14, 2026}
\sourceloc{the six-page author corrections document, physical page 1,
Chapter 1 entry headed ``Subsection 1.10 (Coend), clarification''}
\noindent
\begin{minipage}[t]{0.47\linewidth}\raggedright
\textbf{2015 source, retained above.}
The direct sum in (1.9) is printed with the index $X\in\cC$.
The two limits are printed without an explicit diagram index.
The formulas above reproduce those original expressions.
\end{minipage}\hfill
\begin{minipage}[t]{0.47\linewidth}\raggedright
\textbf{2026 clarification, separate source.}
The authors specify representatives of isomorphism classes and a diagram
$J$ for the (co)limits. Their complete clarification entry is transcribed
immediately below.
\end{minipage}

\medskip
\noindent\textbf{Author clarification (literal transcription).}
\emph{Subsection 1.10} (Coend), clarification. In (1.9), the direct sum is
taken over isomorphism classes of objects. Also the (co)limits in lines 6,7
of the section are taken for the following diagram $J$: the vertices $v_X$
of $J$ are labeled by representatives $X$ of isomorphism classes of objects,
and arrows are $\phi_f:v_X\to v_{X\oplus Y}$ and
$\psi_f:v_Y\to v_{X\oplus Y}$ for every morphism $f:X\to Y$.
The map $E_F:J\to\cC$ is given by the formula
\[
E_F(v_X)=\End(F(X)),\qquad
E_F(\phi_f)(a)=f\circ a,\qquad E_F(\psi_f)(a)=a\circ f,
\]
where we view $\Hom(X,Y)$ naturally as a subgroup of $\End(X\oplus Y)$.
Then $\varinjlim\End(F(X))$ reproduces the usual definition of $\End(F)$:
the set of collections $\{a_X\in\End(F(X))\}$ such that
$a_Y\circ f=f\circ a_Y$ for all $X,Y$ and $f:X\to Y$.

\editorial{This separate entry reproduces the symbols printed in the 2026
document, including $E_F:J\to\cC$, its formulas involving $f$, the
rightward limit arrow, and both occurrences of $a_Y$ in its final equation.
They are not silently emended here. The original theorem and sketch above
retain the 2015 wording. The ind-object $Q$, its quotient $I$, and the
extension implicit in $F(I)$ are not asserted here to be objects of
$\cC$ itself or to be finite dimensional.}
% END SUPPORT AUTHOR CLARIFICATION EGNO:1.10.1

\end{document}
